% J3.6: flujo de trabajo con OrdinaryDiffEq.jl (de la ecuación de movimiento a la solución).
\documentclass[tikz,border=4pt]{standalone}
\usepackage{resonancia}
\begin{document}
\begin{tikzpicture}
  \node[cajaAzul, minimum width=3.0cm, minimum height=1.3cm] (e) at (0,0) {\textbf{Ecuación}\\[1pt]{\footnotesize $\dot{\mathbf u}=\mathbf f(\mathbf u,p,t)$}};
  \node[cajaVerde, minimum width=3.6cm, minimum height=1.3cm] (p) at (5.0,0) {\textbf{Problema}\\[1pt]{\scriptsize \texttt{ODEProblem(f!,u0,tspan,p)}}};
  \node[cajaNaranja, minimum width=3.6cm, minimum height=1.3cm] (s) at (10.3,0) {\textbf{Solución}\\[1pt]{\scriptsize \texttt{solve(prob,alg;reltol,abstol)}}};
  \node[caja, draw=rVioleta, fill=rVioleta!8, minimum width=3.4cm, minimum height=1.3cm] (u) at (15.5,0) {\textbf{Uso}\\[1pt]{\scriptsize \texttt{sol(t)}, \texttt{saveat}, \texttt{sol.stats}}};
  \draw[flecha] (e) -- (p); \draw[flecha] (p) -- (s); \draw[flecha] (s) -- (u);
  \node[etiqueta, color=rTinta, anchor=south] at (7.65,0.85) {algoritmo y tolerancias};
  \node[etiqueta, color=rGris, anchor=north] at (7.65,-1.0) {\texttt{Tsit5} (no rígido) \quad \texttt{Rodas5P}, \texttt{FBDF} (rígido)};
\end{tikzpicture}
\end{document}
